\section{Limitations}

Our experiments cover seven English BEIR collections, two instruction-tuned
generators, two dense encoders, and BM25 as the sparse retriever. The results
therefore do not establish the same behavior for other languages, learned
sparse retrievers, or other fusion rules. Score-product anchoring additionally
assumes a non-negative sparse scoring function, and the access analysis fixes
equally weighted WRRF, $K=20$, and a specific incremental replay policy.

The reported access measurements are logical certification depths obtained by
replaying stored rankings. They are not wall-clock latency and do not include
reference generation, query encoding, indexing, or communication costs. An
online system would need an incremental retrieval interface before a reduction
in replay depth could be interpreted as an end-to-end speedup.

The direct QuDAR comparison is limited to QuDAR-simple RRF. It follows the
official four-signal, Top-1000 rank-fusion configuration and reuses the same
generated evidence, but our frozen ranking artifacts do not contain the
normalized retrieval scores required by QuDAR-confidence, and we do not add
new LLM relevance judgments for QuDAR-llm. The comparison therefore supports
parity with the reconstructable rank-based variant, not with every QuDAR
weighting strategy.

Generated references may contain irrelevant or unsupported expressions despite
the intent-preserving prompt and structured validation. We average three draws
and report parser fallback behavior, but this does not cover the full variation
of possible generators and decoding settings. Finally, the benchmark relevance
judgments are incomplete, as is common in pooled retrieval collections; methods
that retrieve documents outside the original pools may therefore be
underestimated. Our MuGI comparison reproduces its generation and integration
components without pseudo-relevance-feedback calibration, and the complete
prior-method comparison is limited to the four datasets on which all reproduced
pipelines were available.
